\Article{7}
After \textsc{d'Alembert}, the distinguished \textsc{Euler} published his investigations on the same subject in \textit{Recherches sur les racines imaginaires des équations, Hist. de l'Acad. de Berlin A.} 1749, p. 223 ff. He gave two methods there; the substance of the first is as follows.

First the distinguished Euler undertakes to prove that, if $m$ denotes any power of the number 2, the function $x^{2m}+Bx^{2m-2}+Cx^{2m-3}+\text{etc.}+M=X$ (in which the coefficient of the second term is $=0$) can always be resolved into two real factors in which $x$ rises to degree $m$. For this purpose he assumes two factors,
\[
x^m-ux^{m-1}+\alpha x^{m-2}+\beta x^{m-3}+\text{etc.},\quad\text{and}\quad x^m+ux^{m-1}+\lambda x^{m-2}+\mu x^{m-3}\ \text{etc.}
\]
where the coefficients $u,\alpha,\beta$, etc., $\lambda,\mu$, etc. are still unknown, and sets their product equal to the function $X$. Comparing coefficients then gives $2m-1$ equations, and evidently all that must be proved is that real values satisfying these equations can be assigned to the unknowns $u,\alpha,\beta$, etc., $\lambda,\mu$, etc. (whose number is also $2m-1$). Euler now asserts that, if $u$ is first regarded as known, so that the number of unknowns is one less than the number of equations, then by suitably combining the equations according to known algebraic methods, all of $\alpha,\beta$, etc., $\lambda,\mu$, etc. can be determined rationally, without any extraction of roots, from $u$ and the coefficients $B,C$, etc., and thus attain real values as soon as $u$ becomes real. Moreover, all of $\alpha,\beta$, etc., $\lambda,\mu$, etc. can be eliminated to give an equation $U=0$, where $U$ is a polynomial in $u$ alone and the known coefficients. Expanding this equation itself by the usual method of elimination would be an immense task when the proposed equation $X=0$ is of a rather\SourceLeaf{017} high degree, and completely impossible for an unspecified degree (according to Euler himself, p. 239). Yet here it is enough to know one property of this equation: the constant term in $U$ (which does not involve the unknown $u$) is necessarily negative. From this it is known to follow that the equation has at least one real root, or that $u$, and therefore also $\alpha,\beta$, etc., $\lambda,\mu$, etc., can be determined as real quantities in at least one way. That property can be established by the following considerations. Since $x^m-ux^{m-1}+\alpha x^{m-2}+\text{etc.}$ is supposed to be a factor of the function $X$, $u$ must be the sum of $m$ roots of the equation $X=0$, and must therefore have as many values as there are different ways of choosing $m$ out of $2m$ roots; that is, by the principles of combinatorics, $\frac{2m\mathbin{.}\overline{2m-1}\mathbin{.}\overline{2m-2}.\,.\,.\,.\,.\overline{m+1}}{1\mathbin{.}2\mathbin{.}3\ldots m}$ values. This number is always twice an odd number (I omit the proof, which is not difficult); if it is put $=2k$, its half $k$ is therefore odd, while the equation $U=0$ has degree $2k$. Now, since the second term is absent in the equation $X=0$, the sum of all $2m$ roots is 0. Thus, if the sum of any $m$ roots is $+p$, the sum of the remaining ones is $-p$: that is, if $+p$ is among the values of $u$, then $-p$ is among them too. From this Euler concludes that $U$ is the product of $k$ quadratic factors such as $uu-pp$, $uu-qq$, $uu-rr$, etc., where $+p$, $-p$, $+q$, $-q$, etc. denote all $2k$ roots of the equation $U=0$. Because the number of these factors is odd, the constant term in $U$ is therefore the square of the product $pqr$, etc., with a negative sign. But the product $pqr$, etc. can always be determined rationally from the coefficients $B,C$, etc., and must therefore be a real quantity. Its square with a negative sign is thus certainly a negative quantity. Q. E. D.

Since these two real factors of $X$ have degree $m$, and $m$ is a power of the number 2, each can in the same way be resolved again into two real factors of degree $\frac12m$. Repeated halving of the number $m$ must eventually lead to two; hence continuing the operation evidently gives the function $X$ resolved into real quadratic factors.

If, however, a function is proposed in which the second term is not absent, say $x^{2m}+Ax^{2m-1}+Bx^{2m-2}+\text{etc.}+M$, where $2m$ still denotes a power of two, the substitution $x=y-\frac{A}{2m}$ transforms it into a similar function without a second term. It readily follows that this function too can be resolved into real quadratic factors.

Finally, given a function of degree $n$, where $n$ denotes a number that is not
\SourceLeaf{018}
a power of two, let the next larger power of two than $n$ be put $=2m$, and multiply the proposed function by any $2m-n$ real linear factors. From the fact that the product can be resolved into real quadratic factors, it readily follows that the proposed function too must be resolvable into real factors of the second or first degree.
